% GENERATED FILE. Edit canonical lessons, then run npm run generate:books.
% canonical-source-hash: fnv1a64:b065b31386d0106a
% canonical-lessons: KA-C477-tola, KA-C477-kaduhandi, KA-C477-gharjisu, KA-C477-meyu, KA-C477-geddalu

\chapter{Creatures of Forest and Field}
\label{ch:ka-creatures-of-forest-and-field}
\hlchaptermodality{477}

\begin{chapteropening}
\textbf{By the end of this chapter:} \emph{I can say a wolf, a wild boar, to roar, to graze and a termite.}

Complete the last lesson of chapter 477: I can say a wolf, a wild boar, to roar, to graze and a termite.
\end{chapteropening}

\section[\texorpdfstring{\kn{ತೋಳ} (tōḷa)}{ತೋಳ (tōḷa)}]{\kn{ತೋಳ} (tōḷa) — a wolf}

\label{lesson:KA-C477-tola}



\begin{quote}
Before the new one: say the Kannada for biodiversity, then the Kannada for to migrate.
\end{quote}

\subsection*{You'll want to know: \kn{ತೋಳ}}
\textbf{\kn{ತೋಳ}} — \emph{tōḷa} — \textquotedblleft{}a wolf\textquotedblright{}.

\begin{quote}
\kn{ತೋಳಗಳು} \kn{ರಾತ್ರಿ} \kn{ಗುಂಪಿನಲ್ಲಿ} \kn{ತಿರುಗುತ್ತವೆ}.
\end{quote}

\begin{quote}
\emph{tōḷagaḷu rātri gumpinalli tiruguttave.}
\end{quote}

\begin{quote}
Wolves roam about in packs at night.
\end{quote}

Plural: \kn{ತೋಳಗಳು}, \emph{tōḷagaḷu}.

The plural takes -\kn{ಗಳು}: \kn{ತೋಳಗಳು}.

\begin{cousinweb}
\textbf{\kn{ತೋಳ}} (\emph{tōḷa}) is native Dravidian: DEDR 3548 lists Kannada \emph{tōḷa}, \textquotedblleft{}a wolf\textquotedblright{}.
\end{cousinweb}

\begin{grammarlens}[title={What you've built}]
One more word for getting around.
\end{grammarlens}

\subsection*{Your turn}
\begin{itemize}
\raggedright
  \item \emph{Say it:} \emph{tōḷa}
  \item \emph{Say it:} \emph{tōḷa}, once more
  \item \emph{Say it:} \emph{tōḷagaḷu rātri gumpinalli tiruguttave.} — Wolves roam about in packs at night.
  \item \emph{Recall:} say the Kannada for biodiversity, then the Kannada for to migrate, then say \emph{tōḷa} again
\end{itemize}

\begin{tcolorbox}[breakable,colback=teal!4,colframe=teal!35!black,title={Before you move on}]
What does \kn{ತೋಳ} mean? (\textquotedblleft{}A wolf\textquotedblright{}.) Say it once more.
\end{tcolorbox}
\section[\texorpdfstring{\kn{ಕಾಡುಹಂದಿ} (kāḍuhandi)}{ಕಾಡುಹಂದಿ (kāḍuhandi)}]{\kn{ಕಾಡುಹಂದಿ} (kāḍuhandi) — a wild boar}

\label{lesson:KA-C477-kaduhandi}



\begin{quote}
Before the new one: say the Kannada for to migrate, then the Kannada for a wolf.
\end{quote}

\subsection*{You'll want to know: \kn{ಕಾಡುಹಂದಿ}}
\textbf{\kn{ಕಾಡುಹಂದಿ}} — \emph{kāḍuhandi} — \textquotedblleft{}a wild boar\textquotedblright{}.

\begin{quote}
\kn{ಕಾಡುಹಂದಿಗಳು} \kn{ರೈತನ} \kn{ಹೊಲವನ್ನು} \kn{ಹಾಳು} \kn{ಮಾಡಿದವು}.
\end{quote}

\begin{quote}
\emph{kāḍuhandigaḷu raitana holavannu hāḷu māḍidavu.}
\end{quote}

\begin{quote}
Wild boars ruined the farmer's field.
\end{quote}

\kn{ಕಾಡು}, forest, plus \kn{ಹಂದಿ}, a pig.

\begin{cousinweb}
\textbf{\kn{ಕಾಡುಹಂದಿ}} (\emph{kāḍuhandi}) is a native Kannada compound of \kn{ಕಾಡು}, \textquotedblleft{}a forest\textquotedblright{} (DEDR 1438), and \kn{ಹಂದಿ}, \textquotedblleft{}a pig\textquotedblright{} (DEDR 4039).
\end{cousinweb}

\begin{grammarlens}[title={What you've built}]
One more word for getting around.
\end{grammarlens}

\subsection*{Your turn}
\begin{itemize}
\raggedright
  \item \emph{Say it:} \emph{kāḍuhandi}
  \item \emph{Say it:} \emph{kāḍuhandi}, once more
  \item \emph{Say it:} \emph{kāḍuhandigaḷu raitana holavannu hāḷu māḍidavu.} — Wild boars ruined the farmer's field.
  \item \emph{Recall:} say the Kannada for to migrate, then the Kannada for a wolf, then say \emph{kāḍuhandi} again
\end{itemize}

\begin{tcolorbox}[breakable,colback=teal!4,colframe=teal!35!black,title={Before you move on}]
What does \kn{ಕಾಡುಹಂದಿ} mean? (\textquotedblleft{}A wild boar\textquotedblright{}.) Say it once more.
\end{tcolorbox}
\section[\texorpdfstring{\kn{ಘರ್ಜಿಸು} (gharjisu)}{ಘರ್ಜಿಸು (gharjisu)}]{\kn{ಘರ್ಜಿಸು} (gharjisu) — to roar}

\label{lesson:KA-C477-gharjisu}



\begin{quote}
Before the new one: say the Kannada for a wolf, then the Kannada for a wild boar.
\end{quote}

\subsection*{You'll want to know: \kn{ಘರ್ಜಿಸು}}
\textbf{\kn{ಘರ್ಜಿಸು}} — \emph{gharjisu} — \textquotedblleft{}to roar\textquotedblright{}.

\begin{quote}
\kn{ರಾತ್ರಿ} \kn{ಕಾಡಿನಲ್ಲಿ} \kn{ಸಿಂಹ} \kn{ಘರ್ಜಿಸಿತು}.
\end{quote}

\begin{quote}
\emph{rātri kāḍinalli siṃha gharjisitu.}
\end{quote}

\begin{quote}
At night a lion roared in the forest.
\end{quote}

Past (it …): \kn{ಘರ್ಜಿಸಿತು}, \emph{gharjisitu}.

Of lions and tigers, and of thunder.

\begin{cousinweb}
\textbf{\kn{ಘರ್ಜಿಸು}} (\emph{gharjisu}) is built from Sanskrit \emph{garj}, \textquotedblleft{}to thunder, to roar\textquotedblright{} (CDIAL 4046, \emph{garjati}), with the verb ending -\emph{isu}; Kannada writes it with the aspirate \kn{ಘ}.
\end{cousinweb}

\begin{grammarlens}[title={What you've built}]
One more word for getting around.
\end{grammarlens}

\subsection*{Your turn}
\begin{itemize}
\raggedright
  \item \emph{Say it:} \emph{gharjisu}
  \item \emph{Say it:} \emph{gharjisu}, once more
  \item \emph{Say it:} \emph{rātri kāḍinalli siṃha gharjisitu.} — At night a lion roared in the forest.
  \item \emph{Recall:} say the Kannada for a wolf, then the Kannada for a wild boar, then say \emph{gharjisu} again
\end{itemize}

\begin{tcolorbox}[breakable,colback=teal!4,colframe=teal!35!black,title={Before you move on}]
What does \kn{ಘರ್ಜಿಸು} mean? (\textquotedblleft{}To roar\textquotedblright{}.) Say it once more.
\end{tcolorbox}
\section[\texorpdfstring{\kn{ಮೇಯು} (mēyu)}{ಮೇಯು (mēyu)}]{\kn{ಮೇಯು} (mēyu) — to graze}

\label{lesson:KA-C477-meyu}



\begin{quote}
Before the new one: say the Kannada for a wild boar, then the Kannada for to roar.
\end{quote}

\subsection*{You'll want to know: \kn{ಮೇಯು}}
\textbf{\kn{ಮೇಯು}} — \emph{mēyu} — \textquotedblleft{}to graze\textquotedblright{}.

\begin{quote}
\kn{ಹಸುಗಳು} \kn{ಹೊಲದಲ್ಲಿ} \kn{ಹುಲ್ಲು} \kn{ಮೇಯುತ್ತಿದ್ದವು}.
\end{quote}

\begin{quote}
\emph{hasugaḷu holadalli hullu mēyuttiddavu.}
\end{quote}

\begin{quote}
The cows were grazing on grass in the field.
\end{quote}

Past (it …): \kn{ಮೇಯಿತು}, \emph{mēyitu}.

Of cattle and deer; \kn{ಮೇಯಿಸು} is to take animals to graze.

\begin{cousinweb}
\textbf{\kn{ಮೇಯು}} (\emph{mēyu}) is native Dravidian: DEDR 5093 lists Kannada \emph{mē}, \emph{mēyu}, \textquotedblleft{}to graze, to eat grass\textquotedblright{}.
\end{cousinweb}

\begin{grammarlens}[title={What you've built}]
One more word for getting around.
\end{grammarlens}

\subsection*{Your turn}
\begin{itemize}
\raggedright
  \item \emph{Say it:} \emph{mēyu}
  \item \emph{Say it:} \emph{mēyu}, once more
  \item \emph{Say it:} \emph{hasugaḷu holadalli hullu mēyuttiddavu.} — The cows were grazing on grass in the field.
  \item \emph{Recall:} say the Kannada for a wild boar, then the Kannada for to roar, then say \emph{mēyu} again
\end{itemize}

\begin{tcolorbox}[breakable,colback=teal!4,colframe=teal!35!black,title={Before you move on}]
What does \kn{ಮೇಯು} mean? (\textquotedblleft{}To graze\textquotedblright{}.) Say it once more.
\end{tcolorbox}
\section[\texorpdfstring{\kn{ಗೆದ್ದಲು} (geddalu)}{ಗೆದ್ದಲು (geddalu)}]{\kn{ಗೆದ್ದಲು} (geddalu) — a termite}

\label{lesson:KA-C477-geddalu}



\begin{quote}
Before the new one: say the Kannada for to roar, then the Kannada for to graze.
\end{quote}

\subsection*{You'll want to know: \kn{ಗೆದ್ದಲು}}
\textbf{\kn{ಗೆದ್ದಲು}} — \emph{geddalu} — \textquotedblleft{}a termite\textquotedblright{}.

\begin{quote}
\kn{ಗೆದ್ದಲು} \kn{ಹಳೆಯ} \kn{ಬಾಗಿಲನ್ನು} \kn{ತಿಂದುಹಾಕಿತು}.
\end{quote}

\begin{quote}
\emph{geddalu haḷeya bāgilannu tinduhākitu.}
\end{quote}

\begin{quote}
Termites ate away the old door.
\end{quote}

Also called white ants; their mound is a \kn{ಹುತ್ತ}.

\begin{cousinweb}
\textbf{\kn{ಗೆದ್ದಲು}} (\emph{geddalu}) is native Dravidian: DEDR 1548 lists Kannada \emph{gedal}, \emph{geddalu}, \textquotedblleft{}a termite, a white ant\textquotedblright{}.
\end{cousinweb}

\begin{grammarlens}[title={What you've built}]
One more word for getting around.
\end{grammarlens}

\subsection*{Your turn}
\begin{itemize}
\raggedright
  \item \emph{Say it:} \emph{geddalu}
  \item \emph{Say it:} \emph{geddalu}, once more
  \item \emph{Say it:} \emph{geddalu haḷeya bāgilannu tinduhākitu.} — Termites ate away the old door.
  \item \emph{Recall:} say the Kannada for to roar, then the Kannada for to graze, then say \emph{geddalu} again
\end{itemize}

\begin{tcolorbox}[breakable,colback=teal!4,colframe=teal!35!black,title={Before you move on}]
What does \kn{ಗೆದ್ದಲು} mean? (\textquotedblleft{}A termite\textquotedblright{}.) Say it once more.
\end{tcolorbox}
